% ==========================================================================
% GENERATED by scripts/split_paper.py. Do not edit this file.
%
% Source: imported_papers/Koopman-JEPA-TMLR-rewrite/main.tex
% Part 10: The Soft Spectral Penalty
%
% Edits belong upstream, in the published paper. To change what the
% chapter includes, change imports/bounded.toml and re-run
% the splitter.
% ==========================================================================
\section{The Soft Spectral Penalty}
\label{app:soft}
% [orig 4.2] DONE, verbatim. The two equations merge into one display, and the
% forward reference to the results now points at Section 5.1.
This appendix defines the constraint that the Koopman~(Soft) variant of Section~\ref{sec:variants} applies.

The soft variant penalizes eigenvalues of $\K$ outside the unit circle. The bound of Equation~\ref{eq:linear_bound} requires the spectral norm, $\|\K\|_2 \le 1$. We penalize the spectral radius $\rho(\K)$ instead, which governs the asymptotic growth rate $\rho(\K)^k$ but permits transient growth (see the footnote following Equation~\ref{eq:linear_bound}). Unlike the Cayley variant, the soft penalty leaves the model free to place eigenvalues strictly inside the unit circle and so to represent damping. A soft spectral norm penalty, which power iteration makes cheap~\citep{Miyato2018}, would keep that freedom under the stronger condition and is a natural extension of the variant studied here. The penalty is
\begin{equation}
    \mathcal{L}_\text{soft} = \max(0,\; \rho(\K) - 1),
    \qquad
    \mathcal{L} = \mathcal{L}_\text{base} + \lambda_s\,\mathcal{L}_\text{soft},
    \label{eq:soft_penalty}
\end{equation}
so it activates only when $\rho(\K) > 1$. The weight $\lambda_s$ is selected per method (Table~\ref{tab:hyperparams}). Section~\ref{sec:results} reports what this weaker constraint buys.

% =============================================================================
